\input{../tex-inputs/latex-leseansicht-vorspann}

\section[Arthur Schnitzler an Gustav Schwarzkopf, 8. 7. 1907]{L04451 Arthur Schnitzler an Gustav Schwarzkopf, 8. 7. 1907}
\nopagebreak\mylabel{L04451v}
\rehead{ }\normalsize\beginnumbering\briefempfaengerindex{Schwarzkopf, Gustav@\textsc{Schwarzkopf, Gustav}!zzzSchnitzler, Arthur@\emph{von Arthur Schnitzler}!1907-07-081@{8. 7. 1907}|(be}
\toendnotes[C]{\smallbreak\pagebreak[2]}
\correspDesc{Versand  durch Arthur Schnitzler am 8. 7. 1907 in Wildbad Waldbrunn
\newline{}Übermittlung  am 8. 7. 1907 in Wildbad Waldbrunn
\newline{}Übermittlung  am 9. 7. 1907 in Welsberg-Taisten
\newline{}Erhalt  durch Gustav Schwarzkopf im Zeitraum [9. 7. 1907 – 13. 7. 1907?] in Wien}\toendnotes[C]{\smallbreak}
\Standort{DLA, A:Schnitzler, HS.1985.1.1897.}
\physDesc{Bildpostkarte, 424 Zeichen
\newline{}Handschrift: Bleistift, deutsche Kurrent
\newline{}Versand: 1) Stempel: »\nobreak{}\oindex{Wildbad Waldbrunn@\textbf{Wildbad Waldbrunn}|pwk}Wildbad Waldbrunn, 8. Jul. 1907\nobreak{}«.   2) Stempel: »\nobreak{}\oindex{Welsberg-Taisten@\textbf{Welsberg-Taisten}|pwk}W{[}elsberg{]}, \textcolor{gray}{9. 7. 07}\nobreak{}«. }\toendnotes[C]{\smallbreak}\pstart{}{\pb}Herrn \textsc{Gustav
                     Schwarzkopf}\pend{}\pstart{}Wien I\oindex{I., Innere Stadt@\textbf{I., Innere Stadt}|pw}\pend{}\pstart{}\textsc{Tiefer
                        Graben 23}\oindex{Wien@\textbf{Wien}!I., Innere Stadt@\textbf{I., Innere Stadt}!Tiefer Graben 23@\textbf{Tiefer Graben 23}|pw}\pend{}{\bigskip}
\pstart
           \noindent{}\centering{}{\pb}\textcolor{gray}{\textbf{Bad Waldbrunn\oindex{Wildbad Waldbrunn@\textbf{Wildbad Waldbrunn}|pw} bei Welsberg\oindex{Welsberg-Taisten@\textbf{Welsberg-Taisten}|pw}}}\pend
           \vspace{1em}
\pstart
           \noindent{}{\pb}\label{T_L04451-1v}\edtext{Hier, lieber Guſtav,{ }\label{K_L04451-1v}\edtext{wie vor
               6 Jahren}{\lemma{\textnormal{\emph{wie vor
               6 Jahren}}}\Cendnote{\textnormal{Arthur und Olga
                  Schnitzler\pwindex{Schnitzler, Olga 17.\,1.\,1882 Wien – 13.\,1.\,1970 Lugano@\textsc{Schnitzler, Olga} (17.\,1.\,1882 Wien – 13.\,1.\,1970 Lugano)|pwk} hielten sich vom 15. 8. 1901 bis zum 26. 8. 1901 im Wildbad Waldbrunn\oindex{Wildbad Waldbrunn@\textbf{Wildbad Waldbrunn}|pwk}
                  bei Welsberg\oindex{Welsberg-Taisten@\textbf{Welsberg-Taisten}|pwk} auf.}}}\label{K_L04451-1}, haben wir uns
               angeſiedelt und ſind (vom Wetter abgeſehn) von Gegend und Hotel\oindex{Wildbad Waldbrunn@\textbf{Wildbad Waldbrunn}|pwv} wieder ſehr angethan (um nicht zu{ }ſagen, entzückt)}{\lemma{\textnormal{\emph{Hier, … entzückt)}}}\Cendnote{\textnormal{Der erste Absatz ist unter das Bildmotiv geschrieben.}}}\label{T_L04451-1}\pend
           
\pstart
           \label{T_L04451-2v}\edtext{\label{K_L04451-2v}\edtext{Der Bub\pwindex{Schnitzler, Heinrich 9.\,8.\,1902 Hinterbrühl – 12.\,7.\,1982 Wien@\textsc{Schnitzler, Heinrich} (9.\,8.\,1902 Hinterbrühl – 12.\,7.\,1982 Wien)|pwv} ko{\geminationm}t nach}{\lemma{\textnormal{\emph{Der Bub kommt nach}}}\Cendnote{\textnormal{Am 5. 7. 1907 notierte Schnitzler im \emph{Tagebuch}\pwindex{Schnitzler, Arthur 15.\,5.\,1862 Wien – 21.\,10.\,1931 ebd.@\textsc{Schnitzler, Arthur} (15.\,5.\,1862 Wien – 21.\,10.\,1931 ebd.), \emph{Schriftsteller, Mediziner}!Tagebuch@\strich\emph{Tagebuch}|pwk}
                     den »Entschluss hierzubleiben und den Buben\pwindex{Schnitzler, Heinrich 9.\,8.\,1902 Hinterbrühl – 12.\,7.\,1982 Wien@\textsc{Schnitzler, Heinrich} (9.\,8.\,1902 Hinterbrühl – 12.\,7.\,1982 Wien)|pwv} nachkommen zu lassen.« Am 10. 7. 1907 vermerkte
                     er »Heini’s\pwindex{Schnitzler, Heinrich 9.\,8.\,1902 Hinterbrühl – 12.\,7.\,1982 Wien@\textsc{Schnitzler, Heinrich} (9.\,8.\,1902 Hinterbrühl – 12.\,7.\,1982 Wien)|pw} Ankunft«.}}}\label{K_L04451-2} und
               wir bleiben hier. Von Sekirn\oindex{Sekirn@\textbf{Sekirn}|pw}{ }ſind ein paar
               Stunden her, we{\geminationn} Sie überhaupt hinfahren  müſſen Sie
               unbedingt}{\lemma{\textnormal{\emph{Der … unbedingt}}}\Cendnote{\textnormal{am rechten Rand der Karte geschrieben}}}\label{T_L04451-2}{ }\label{T_L04451-3v}\edtext{auch hieher ko{\geminationm}en. Jedenfalls also laſſen Sie
               ein Wort von{ }ſich hören.}{\lemma{\textnormal{\emph{auch … hören.}}}\Cendnote{\textnormal{Text über dem Bild um 180° gedreht}}}\label{T_L04451-3}\pend
           
\pstart
           Wir grüßen {\\[\baselineskip]}herzlichſt{\\[\baselineskip]}\spacefill\mbox{A.}\pend
           \leftskip=0em{}\selectlanguage{ngerman}\endnumbering\briefempfaengerindex{Schwarzkopf, Gustav@\textsc{Schwarzkopf, Gustav}!zzzSchnitzler, Arthur@\emph{von Arthur Schnitzler}!1907-07-081@{8. 7. 1907}|)be}\mylabel{L04451h}
\begin{anhang}
\end{anhang}\newcommand{\dateiname}{L04451}\newcommand{\titel}{Arthur Schnitzler an Gustav Schwarzkopf, 8. 7. 1907}\newcommand{\editorInnen}{Herausgegeben von Jahnke, SelmaMüller, Martin Anton}\input{../tex-inputs/latex-leseansicht-abspann}
